\documentclass[12pt,a4paper]{article}

\usepackage[utf8]{inputenc}
\usepackage[T1]{fontenc}
\usepackage[brazil]{babel}
\usepackage{geometry}
\usepackage{hyperref}
\usepackage{setspace}
\usepackage{enumitem}
\usepackage{graphicx}
\usepackage{titlesec}
\usepackage{float}

\geometry{a4paper, margin=2.5cm}
\setstretch{1.3}
\hypersetup{
    colorlinks=true,
    linkcolor=black,
    urlcolor=blue,
    citecolor=blue,
    bookmarksopen=true,
    pdfstartview=FitH
}

\setlist{itemsep=0.5em, parsep=0pt, topsep=0.5em}
\setlist[enumerate,1]{itemsep=0.8em}
\setlist[itemize,1]{itemsep=0.5em}

\titlespacing*{\section}{0pt}{3ex plus 1ex minus .2ex}{2.3ex plus .2ex}
\titlespacing*{\subsection}{0pt}{2.5ex plus 1ex minus .2ex}{1.5ex plus .2ex}
\titlespacing*{\subsubsection}{0pt}{2ex plus 1ex minus .2ex}{1ex plus .2ex}

\begin{document}

\title{
    \textbf{Relatório Final do Trabalho de Programação Paralela} \\[0.5em]
    \large Cálculo de Números Primos com MPI \\[0.3em]
    \normalsize Comparação entre Naive e Bag of Tasks
}
\author{
    Eduardo Rottschaefer Oliveira \\
    Breno Alvarenga \\
    Marcus Vinicius Bossan \\
}
\date{\today}

\maketitle
\thispagestyle{empty}

\newpage
\tableofcontents
\newpage

\section{Introdução}
Este trabalho apresenta a implementação e a avaliação de duas estratégias de paralelização para o cálculo de números primos com MPI: uma versão \textit{naive}, com divisão estática do intervalo, e uma versão \textit{bag of tasks}, com distribuição dinâmica de blocos de trabalho.

O objetivo foi comparar diferentes combinações de comunicação ponto a ponto usando \texttt{MPI\_Send}, \texttt{MPI\_Isend}, \texttt{MPI\_Rsend}, \texttt{MPI\_Bsend}, \texttt{MPI\_Ssend}, combinadas com \texttt{MPI\_Recv} ou \texttt{MPI\_Irecv}, executando os testes com 4 processos e limite de entrada igual a 500.000.000.

\section{Desenvolvimento}

Na implementação \textit{naive}, cada processo calcula uma faixa fixa do intervalo e envia a contagem parcial ao processo raiz. Na implementação \textit{bag of tasks}, o processo raiz distribui blocos de trabalho conforme recebe resultados, o que tende a reduzir o tempo ocioso dos processos trabalhadores.

\section{Metodologia}
Os testes foram executados com 4 processos MPI e limite de entrada igual a 500.000.000. O ambiente de execução utilizado conta com um processador 11\textsuperscript{a} Geração Intel\textregistered\ Core\texttrademark\ i3-1115G4, composto por 2 núcleos físicos e suporte a \textit{Hyper-Threading} (processador visto como possuindo 4 lógicos/threads pelo sistema operacional). O tempo de execução foi medido pelo processo raiz ao final da execução, permitindo comparar o custo das diferentes estratégias de comunicação.

\section{Resultados}
\subsection{Naive}

\subsubsection{Comunicação com \texttt{MPI\_Recv}}
\begin{table}[H]
\centering
\begin{tabular}{|c|c|c|c|}
\hline
\textbf{Modo de envio} & \textbf{Tempo (s)} & \textbf{Speedup} & \textbf{Eficiência} \\
\hline
MPI\_Send & 488.121 & 1,386 & 0,346 \\
\hline
MPI\_Isend & 472.150 & 1,433 & 0,358 \\
\hline
MPI\_Rsend & 474.950 & 1,425 & 0,356 \\
\hline
MPI\_Bsend & 473.966 & 1,428 & 0,357 \\
\hline
MPI\_Ssend & 460.119 & 1,470 & 0,367 \\
\hline
\end{tabular}
\caption{Resultados naive utilizando \texttt{MPI\_Recv}.}
\end{table}

\subsubsection{Comunicação com \texttt{MPI\_Irecv}}
\begin{table}[H]
\centering
\begin{tabular}{|c|c|c|c|}
\hline
\textbf{Modo de envio} & \textbf{Tempo (s)} & \textbf{Speedup} & \textbf{Eficiência} \\
\hline
MPI\_Send & 452.886 & 1,494 & 0,374 \\
\hline
MPI\_Isend & 428.942 & 1,578 & 0,394 \\
\hline
MPI\_Rsend & 415.228 & 1,630 & 0,408 \\
\hline
MPI\_Bsend & 455.772 & 1,485 & 0,371 \\
\hline
MPI\_Ssend & 477.354 & 1,418 & 0,354 \\
\hline
\end{tabular}
\caption{Resultados naive utilizando \texttt{MPI\_Irecv}.}
\end{table}

Na versão naive com \texttt{MPI\_Recv}, o menor tempo foi obtido com \texttt{MPI\_Ssend}, enquanto o maior tempo ocorreu com \texttt{MPI\_Send}. Já na versão com \texttt{MPI\_Irecv}, o menor tempo foi obtido com \texttt{MPI\_Rsend}, enquanto o maior tempo ocorreu com \texttt{MPI\_Ssend}.

\subsection{Bag of Tasks}

\subsubsection{Comunicação com \texttt{MPI\_Recv}}
\begin{table}[H]
\centering
\begin{tabular}{|c|c|c|c|}
\hline
\textbf{Modo de envio} & \textbf{Tempo (s)} & \textbf{Speedup} & \textbf{Eficiência} \\
\hline
MPI\_Send & 376.948 & 1,796 & 0,449 \\
\hline
MPI\_Isend & 382.908 & 1,768 & 0,442 \\
\hline
MPI\_Bsend & 376.629 & 1,797 & 0,449 \\
\hline
MPI\_Ssend & 375.820 & 1,801 & 0,450 \\
\hline
\end{tabular}
\caption{Resultados bag of tasks utilizando \texttt{MPI\_Recv}.}
\end{table}

\subsubsection{Comunicação com \texttt{MPI\_Irecv}}
\begin{table}[H]
\centering
\begin{tabular}{|c|c|c|c|}
\hline
\textbf{Modo de envio} & \textbf{Tempo (s)} & \textbf{Speedup} & \textbf{Eficiência} \\
\hline
MPI\_Send & 384.107 & 1,762 & 0,441 \\
\hline
MPI\_Isend & 396.906 & 1,705 & 0,426 \\
\hline
MPI\_Bsend & 404.546 & 1,673 & 0,418 \\
\hline
MPI\_Ssend & 435.806 & 1,553 & 0,388 \\
\hline
\end{tabular}
\caption{Resultados bag of tasks utilizando \texttt{MPI\_Irecv}.}
\end{table}

Na versão bag of tasks com \texttt{MPI\_Recv}, o menor tempo foi obtido com \texttt{MPI\_Ssend}, e o maior tempo com \texttt{MPI\_Isend}. Em contrapartida, com o uso de \texttt{MPI\_Irecv}, o menor tempo ocorreu utilizando \texttt{MPI\_Send}, e o pior desempenho foi com \texttt{MPI\_Ssend}.

\section{Conclusão}
A versão bag of tasks apresentou desempenho superior à versão naive em todas as combinações registradas. Isso é compatível com o comportamento esperado de uma distribuição dinâmica de carga, que reduz o tempo ocioso e melhora o aproveitamento dos processos paralelos.

Considerando as medições disponíveis, a redução de tempo varia de forma consistente entre as combinações testadas, com destaque para o uso de \texttt{MPI\_Ssend}. O conjunto obtido foi suficiente para consolidar a comparação entre as duas estratégias e para demonstrar o ganho prático da abordagem bag of tasks.

\end{document}